\section{The source short reading and its common scalar}
\label{sec:sr-source-short-reading-normalization}

The source comparison is a comparison of the actual framed
section and the actual lattice theta series.  Its normalizing
factor is selected in two stages.  A lattice factor is fixed
before the component; a second factor is fixed on a component
before the lattice, polynomial, short rule or frame is varied.
We retain that order in the statements below.

\subsection{The value formula on one analytic fibre}
Fix the original host \(H\), datum \(D_w\), level \(l\),
moduli object \(X_o\), coefficient field \(F=c.F\) embedded
in \(\mathbb C\), and the actual connection \(\mathrm{nab}\)
with \(\mathrm{IsNab}\).
Let \(X\) be the original Hermitian data of rank \(r\),
with symplectic basis \(\mathrm{bas}\), coordinate index
\(\mathrm{idx}\), and \(x\in X.\mathrm{dom}\).
Let \(L\) be self-dual for the alternating form of \(X\),
and let
\[
 A:\mathrm{AnFib}(X,\mathrm{bas},\mathrm{idx},e,x,L,y)
\]
be the actual analytic fibre of the framed point \(y\).
Thus its exponential, integral period kernel, differential,
\(O_E\)-action and local line trivializations are tied to
that same \(y\).  We use its section-reading relation
\(A.\mathrm{Reads}\), rather than identifying two line bundles
by their ranks.

Let \(n\) be even, \(r=D_w.p+D_w.q\), and let \(e_L\) be
the actual letter bijection satisfying
\[
 \mathrm{letterDir}(e_L(\tau,i,b))
        =e_b\otimes e_{e^{-1}(\tau,i)} .
 \tag{RS1}\label{eq:sr-source-reading-directions}
\]
Take a level-one reading \(\theta\) with
\(A.\mathrm{Reads}(1,y.s,\theta)\).
Let \(v\) be an actual normalized lattice vacuum for
\(O_E^n\otimes L\), with the host semicharacter
\(w_0+\bar w_0=1\), and suppose
\[
       \prod_{b=1}^{n}\theta(z_b)
            =\kappa_0\,\mathrm{thA}(v,x;z).
 \tag{RS2}\label{eq:sr-source-level-one-product-reading}
\]
For \(n=0\), the products have value \(1\); no extraction
of an \(n\)-th root is used.

Fix a marking \(\mathrm{mk}\) with \(A.\mathrm{Marked}(\mathrm{mk})\).
Let \(\Lambda\in\mathrm{Vset}(E,n,l.M)\) and let
\(\mathrm{NOK}(\Lambda,N,l.M)\) hold, including oddness
and positivity of \(N\).  Put
\[
 z=z_\Lambda(N,r)=\Lambda.\mathrm{zCard}(N,r),\quad
 m=m_\Lambda(N,r)=\Lambda.\mathrm{mIdx}(N,r),\quad
 Q_l=\left|\mathrm{LevQ}(D_w,l)^n\right|.
\]
These are the actual source indices.  Let \(Y_q\) be
the actual roots of the marking, satisfying
\(\mathrm{IsRoots}(l,L,\mathrm{mk},N,Y)\).
For a finite level weight \(w_q\), set
\[
 {\cal G}=\sum_q w_q\,\mathrm{heisG}(Y_q)
             \mathrm{projA}_{\Lambda\otimes L}v .
 \tag{RS3}\label{eq:sr-source-reading-gphi}
\]
All sums here are finite.
For \(P\in\mathrm{HHa}(D_w,a,n)\), a short value \(v_0\)
with the source parameter \(\beta\) satisfies
\[
 v_0=\beta^{-1}m^{-1}Q_l^{-1}z\,\kappa_0\,
     \mathrm{thP}\!\left(
       {\cal G},\,
       \mathrm{rename}_{e_L}(\sigma P),x\right).
 \tag{RS4}\label{eq:sr-source-short-value-normalization}
\]
The scalar factors are embedded into \(\mathbb C\) using
the specified coefficient embedding.

We now prove the calculation in (RS4), stating the precise
geometric identities it uses.
The actual even-representative construction supplies
\(u_1,\ldots,u_z\in\Lambda\otimes L\), distinct modulo
\(N(O_E^n\otimes L)\), with
\[
 N u_{i,b}\in2L,\qquad
 \sum_i\mathrm{heisG}(u_i)v
                    =z\,\mathrm{projA}_{\Lambda\otimes L}v,
       \qquad z,m>0.
 \tag{RS5}\label{eq:sr-source-even-representatives}
\]
To see why doubled representatives exist, if an odd
integer \(N^2\) annihilates a finite quotient, multiplication
by \(2\) is bijective on it: \(2((N^2+1)/2)\equiv1\bmod N^2\).
Choose representatives using this inverse and double them.
The tensor indices are the powers appearing in \(z\).
The signed Heisenberg average is the lattice projector;
the sign of a doubled lattice vector is \(1\), so its
sum is the unsigned sum in (RS5).
This is the same finite-group projection argument as in
Section~\ref{sec:sr-source-raw-theta-norm}.

The analytic-fibre root and isotropic-point identities
turn \(N^{-1}u_i\) and \(N^{-1}Y_q\) into the geometric
torsion points of \(\mathrm{ShortValAt}\).
For each \(q,i,b\), choose its canonical translated
section \(F_{q,i,b}\).
The required actual symmetric-lift reading is
\[
 A.\mathrm{Reads}\left(N^2,F_{q,i,b},
   z_b\longmapsto
   \mathcal F(Y_{q,b})\mathcal F(u_{i,b})\theta(Nz_b)\right),
 \tag{RS6}\label{eq:sr-source-canonical-translation-reading}
\]
where \(\mathcal F\) is the source Fock translation.
Here the right side denotes the two translations before
evaluation at \(Nz_b\).
Their order agrees with
\(\mathrm{heisG}(Y_q)\mathrm{heisG}(u_i)v\).
Oddness is used in the symmetric-lift identity; no
arbitrary lift with the same underlying torsion point
can replace it.

The second required identity is the actual harmonic
connection calculation
\[
 \sum_\mu P_\mu\prod_b
          \mathrm{nab}_y(N^2)(F_{q,i,b},\mathrm{colOf}(\mu,b))
   =P(\partial)\left[
       \kappa_0\,
       \mathrm{thA}(\mathrm{heisG}(Y_q)
                       \mathrm{heisG}(u_i)v,x;Nz)\right]_{z=0},
 \tag{RS7}\label{eq:sr-source-reading-harmonic-connection}
\]
with precisely the directions (RS1).
It is the specialization of the canonical-connection
harmonic contraction to the analytic chart of \(A\).
The formal Wick identity at a supplied chart is one
ingredient of (RS7); the analytic chart's actual germ
and its compatibility with \(\mathrm{IsNab}\) are also
required.

The value definition on this torsion witness is
\[
 \beta^{-1}m^{-1}Q_l^{-1}N^{-\deg P}
    \sum_q w_q\sum_i\sum_\mu P_\mu
             \prod_b\mathrm{nab}_y(N^2)
                           (F_{q,i,b},\mathrm{colOf}(\mu,b)).
\]
It equals the original short value by uniqueness of
\(\mathrm{ShortValAt}\), not by a choice of an unrelated
torsion witness.
Row homogeneity makes \(P\) homogeneous of degree
\(d=\mathrm{DF}(D_w,a)\) when \(P\ne0\).
Thus \(P(\partial)[f(Nz)]_{z=0}=N^dP(\partial)f(0)\);
if \(P=0\), every term is zero, irrespective of the
convention for its total degree.
Use (RS7), commute finite sums and derivatives,
and apply (RS5).  This supplies exactly
\(N^d z\kappa_0\mathrm{thP}({\cal G},\sigma P,x)\).
Since \(N>0\), its \(N^d\) cancels the definition's
\(N^{-\deg P}\), proving (RS4).
Renaming letters commutes with polynomial differentiation
because it bijects the multisets of derivative directions.

\subsection{The geometric Fourier scalar and the lattice factor}
Retain the source cover \(d:\mathrm{DeltaDat}(X_o,c,n)\),
its field identification \(d.e_F\), its actual
\(\mathrm{IsLev}\) witness, and the pinned data
\((\mathrm{ph},\psi,x_0^{\rm pin})\) satisfying
\(d.\mathrm{PhiPin}\).
Use the common classical level \(D\) and the full
\(\mathrm{LevTie}(d,\sigma,\mathrm{ph},D,r\le n)\).
The two level coefficient fields embed compatibly:
\(\sigma_2(d.e_F(t))=\sigma(t)\).
Fix once and for all actual normalized vacua \(v_0(L)\)
for the shifted lattices of \(D\).

The finite Fourier computation of the source has the
following exact shape.  Write
\[
 {\cal C}_\Lambda=
 \frac{\Lambda\otimes(2\bar{\mathfrak m}L)}
      {\Lambda\otimes(2\bar{\mathfrak m}\mathfrak m L)},
       \qquad\mathfrak m=D.\mathfrak m,
 \quad
 \langle u,Y\rangle=\sum_b EF_X(u_b,Y_b).
\]
It first gives a unit scalar \(t_1\), independent of
\(\Lambda,N,Y,q\), with
\[
 \sigma_2(d.w(x_0)(q))
   =t_1\sum_{c\in{\cal C}_\Lambda}
       \mathrm{datF}(D,M_{\rm fr},c)\,
                   e^{-2\pi i\langle c,Y_q\rangle}.
 \tag{RS8}\label{eq:sr-source-cover-fourier-reading}
\]
It also gives the exact finite Heisenberg identity
\[
 \sum_q e^{-2\pi i\langle c,Y_q\rangle}
                         \mathrm{heisG}(Y_q)p
      =Q_{l_2}\,\mathrm{heisG}(c)\,
           \mathrm{projA}_{\Lambda\otimes L^{\rm sh}}p ,
 \tag{RS9}\label{eq:sr-source-roots-finite-fourier}
\]
for \(p\) fixed by the signed lattice action of
\(\Lambda\otimes L\) and \(c\in2\Lambda\otimes\bar{\mathfrak m}L\).
The quotient, phases, shifted lattice and common
Heisenberg representation in these equations are part
of the actual level dictionary.
Equations (RS8) and (RS9) are the explicit lower identities
required here; a separate scalar-existence statement
does not replace them.

The shifted-vacuum projector calculation chooses a
positive \(\rho_{\rm sh}\) and a real function
\(s(\Lambda)\) before \(L,\mathrm{mk},\delta\) or the
vacuum \(v\) is varied.  For each admissible \(L,v\),
its scalar \(t_2\) has
\(|t_2|=\rho_{\rm sh}\), and \(s(\Lambda)>0\)
on the actual \(\mathrm{Vset}(E,n,l_2.M)\).
Substitute (RS8) into the definition (RS3), exchange
the two finite sums, use (RS9), and then this
shifted-vacuum projector identity.  The result is
\[
 {\cal G}=
       t\,s(\Lambda)\,
        \phi_F(D,M_{\rm fr},v_0,L,\Lambda),
 \qquad
 t=t_1Q_{l_2}t_2,\qquad
 |t|=\rho_1:=Q_{l_2}\rho_{\rm sh}>0.
 \tag{RS10}\label{eq:sr-source-gphi-phiF-scalar}
\]
Finiteness of \({\cal C}_\Lambda\) follows from finite
generation and an integer annihilating the quotient,
since \(l_2.M\in\mathfrak m\).
Every passage is a finite sum; there is no convergence
interchange in (RS8)--(RS10).

In the source \(\mathrm{shortRule\_value}\), define,
before \(L,\mathrm{mk},\delta\) are chosen,
\[
 u(\Lambda,N)=
  \sigma_2\!\left(
      d.e_F(\beta_H(c,\Lambda))^{-1}
                   m_\Lambda(N,r)^{-1}Q_{l_2}^{-1}
                   z_\Lambda(N,r)\right)s_+(\Lambda),
 \quad
 s_+(\Lambda)=
  \begin{cases}s(\Lambda),&s(\Lambda)>0,\\1,&\text{otherwise.}\end{cases}
 \tag{RS11}\label{eq:sr-source-uniform-lattice-reading-factor}
\]
The function \(s_+\) simply defines the factor outside
the admissible set as well.
For admissible \(\Lambda,N\), \(u(\Lambda,N)\ne0\).
The indices \(m,z,Q_{l_2}\) are positive.
The source \(\beta_H\) is also nonzero: its definition
selects a nonzero witness when such a witness exists
and otherwise takes \(1\).
This nonvanishing assertion does not require a
covolume or absolute-value assertion about \(\beta_H\).

For each admissible \(L\) and the actual marking/dual/
adaptation ties, the normalized level-one theta factor
has a norm \(\rho_0>0\) independent of its unit frame.
This is the literal mean-one identity
\[
 \frac{\int_{\mathrm{Fd}}|\theta(z)|^2
                      e^{-\pi H_x(z,z)}\,dz}
              {\operatorname{vol}(\mathrm{Fd})}=1
\]
together with (RS2).
It gives
\[
 |\kappa_0|=\rho_0=
 \frac{\left(\sqrt{2^g/\prod_{\tau\in\Psi}X.c_\tau^{\,r}}\right)^{n/2}}
      {|C_{\rm vac}|},
 \tag{RS12}\label{eq:sr-source-mean-one-reading-factor}
\]
where \(C_{\rm vac}\ne0\) is the same period-coordinate
vacuum scalar.  In particular the formula includes
\(n=0\).  Its input includes the actual period-coordinate
theta expansion; equality of one-dimensional spaces
alone would not determine (RS12).

Lift the analytic fibre and marking to the actual
cover, using the given \(\mathrm{IsLev}\) lift, and
apply (RS4) there.
The short-rule definition gives
\(f(y)=d.c_\delta(y_2)v_0\).
By (RS10) and linearity of \(\mathrm{thP}\) in its
test function, the common factor at that fibre is
\[
 \kappa=d.c_\delta(y_2)\,\kappa_0\,t .
\]
The \(\mathrm{UnitC}(d)\) input says
\(|d.c_\delta(y_2)|=1\).
Thus \(|\kappa|=\rho_0\rho_1>0\), independently of
\(\Lambda,N,P,f\), and
\[
 f(y)=\kappa\,u(\Lambda,N)\,
   \mathrm{thP}\!\left(
      \phi_F(D,M_{\rm fr},v_0,L,\Lambda),\,
             \mathrm{rename}_{e_L}(\sigma P),x\right)
 \tag{RS13}\label{eq:sr-source-common-fibre-short-reading}
\]
simultaneously for every one of those choices.
This explains exactly why the scalar can be chosen
before the lattice and polynomial.

\subsection{The scalar is constant on the frame domain}
We prove the analytic step used in the source, including
its simultaneous-family and zero-family cases.

\begin{srlemma}\label{lem:sr-source-common-constant-norm-factor}
Let \(U\) be a connected open subset of a complex normed
vector space, let \(R_i,T_i:U\to\mathbb C\) be holomorphic
for an arbitrary index set \(I\), and let \(\rho>0\).
Suppose that for every \(x\in U\) there is one scalar
\(k_x\), with \(|k_x|=\rho\), such that
\(R_i(x)=k_xT_i(x)\) for every \(i\).
There is \(C\ne0\) with \(R_i=CT_i\) on \(U\) for
every \(i\).
\end{srlemma}
\begin{proof}
If every \(T_i\) is zero, the pointwise identity makes
every \(R_i\) zero; take \(C=\rho\).
Otherwise choose \(T_j(x_0)\ne0\) and a ball about
\(x_0\) on which \(T_j\ne0\).
The holomorphic ratio \(R_j/T_j\) has constant modulus
\(\rho\), so the maximum-modulus principle makes it
constant on that ball, with value \(C\ne0\).
The common pointwise scalar then gives
\(R_i=CT_i\) there for all \(i\).

Here is an explicit continuation argument which also
works in a complex normed vector space.
Let \(S\) be the open set of points at which all these
identities hold in a neighborhood.
It is nonempty.
If \(y\in\overline S\cap U\), choose
\(B(y,\delta)\subset U\) and
\(x_1\in S\cap B(y,\delta/4)\).
For \(z\in B(y,\delta/4)\), the complex line
\(x_1+t(z-x_1)\), \(|t|<3/2\), lies in
\(B(y,\delta)\).
For each \(i\), the holomorphic difference
\(R_i-CT_i\) on that line vanishes near \(t=0\).
The one-variable identity theorem makes it zero
at \(t=1\).
Thus all the identities hold on \(B(y,\delta/4)\);
so \(y\in S\).
The set \(S\) is both open and relatively closed in
the connected \(U\), and hence is \(U\).
\end{proof}

For the source frame version, \(U=X.\mathrm{dom}\)
is open and connected.
The required extension from unit frames is exact:
Gram--Schmidt supplies, for every \(x\in U\), an
upper triangular \(b\) with nonzero diagonal such
that \(xb\) is a unit frame.
The two functions in the reading comparison transform
by the same \(\mathrm{bchar}(\lambda,b)\ne0\).
Cancel that character from their equality at \(xb\);
the same pointwise scalar then works at \(x\).
The preceding lemma therefore applies.
The openness follows from strict leading-minor
inequalities; the connectedness and triangular
orthonormalization are the original frame-domain
results, with the fixed signature of \(X\).

Apply it, for a fixed component \(k\), to the index
set of all admissible quadruples
\((\Lambda,N,P,f)\).
Let \(g_k\) and the coordinate identification
\(\mathrm{castFr}\) be the actual component move
and same-space identification.
Take
\[
 \begin{split}
 R_{\Lambda,N,P,f}(x)&=\mathrm{rdU}_U(f)(k,x),\\
 T_{\Lambda,N,P,f}(x)&=u(\Lambda,N)\,
    \mathrm{thP}\!\left(\phi_F(D,M_{\rm fr},v_0,
                           g_k\Lambda_k,\Lambda),
       \mathrm{rename}_{e_L}(\sigma P),
                    g_k\,\mathrm{castFr}(x)\right).
 \end{split}
\]
The actual section reading and the theta holomorphy
law give holomorphy and the same Borel character.
The common analytic fibre supplied by \(UW\), and
its transported marking, give (RS13) on every unit
frame, with the same \(\rho_0\rho_1\).
Consequently there is \(C_k\ne0\) with
\[
  \mathrm{rdU}_U(f)(k,x)=C_k\,u(\Lambda,N)\,
    \mathrm{thP}\!\left(\phi_F(D,M_{\rm fr},v_0,
                           g_k\Lambda_k,\Lambda),
       \mathrm{rename}_{e_L}(\sigma P),
                    g_k\,\mathrm{castFr}(x)\right).
 \tag{RS14}\label{eq:sr-source-short-reading-global-scalar}
\]
The quantifiers are: fix \(u(\Lambda,N)\) first;
then for each fixed component and its specified
move/marking, choose \(C_k\); then (RS14) holds for
every admissible \(\Lambda,N,P,f\) and every
\(x\in X_k.\mathrm{dom}\).
There is no choice of \(C_k\) depending on \(P\)
or on the individual frame.

\paragraph{Frozen-source correspondence and supplier scope.}
With prefix
\srcid{src/Development/AbelianVarietyTheory/Proofs/SiegelThetaTransformation/},
the value calculation is
\srcid{ShortThetaValueComplexPointEqClassicalTheta.lean}
(\srcid{2055ba448bc5ae903a96758e651e9924a44c367a}),
\srcid{shortValAt_eq_thP};
the common lattice and fibre factors are
\srcid{GphiEqSmulPhiF.lean}
(\srcid{49ecf15b9a7d4565663d55eb915f47c052997523})
and \srcid{IteratedFDerivSmulAll.lean}
(\srcid{b8a0c23b041516b9f4dc141eed167e25c3a52504}),
\srcid{ShortPt.shortRule_value};
the last analytic argument is
\srcid{ConstOfNormFactorUnitFr.lean}
(\srcid{ae99b4fcc02fcee1d1bc16f67e4e521203578ba7})
and \srcid{ShortSeriesReadingsEqClassicalThetaSeries.lean}
(\srcid{569e5e7d4b19f957cfd305524b5d6f1207b233ad}),
\srcid{rdU_short_eq_thP}.
The explicit lower formulas (RS8), (RS9), and (RS12) are
respectively
\srcid{CoverFourierCoeffEqFiniteFourierSum.lean}
(\srcid{74576b27f7c3dd7a54f9c4d66edcd5931e6c9cb4}),
\srcid{FourierTranslatesVacuumRootsLevelPoints.lean}
(\srcid{28fc115f8c46ce96e6f2bbf704fe6c4b1b43aced}),
and \srcid{ExistsSymplOfSelfDual.lean}
(\srcid{8dd7dbad6195af626d9c794fa32465204b5ed6a3}).
The even-representative producer is
\srcid{ExistsEvenRepsTensorLatticeAndSum.lean}
(\srcid{502dfc50708b8b3b17831e6f95df62e543aae8cb}).

The calculation in this section is bounded by its displayed
actual-fibre and finite-Fourier inputs.
Their complete geometric producers are kept separate:
the common \(UW/\mathrm{AnFib}\) realization and period
vacuum expansion; the canonical symmetric-lift reading
(RS6); the analytic-germ and chart-independence
compatibility in (RS7); and the exact Fourier
coefficient/roots identities (RS8)--(RS9).
In particular, the formal Wick contraction alone has
not been identified with the full actual
\(\mathrm{IsNab}\) endpoint.
The exact latter supplier is
\srcid{src/Development/AbelianVarietyTheory/Proofs/CuspFormEvaluations/HarmonicValueCanonicalConnectionEqPolyDeriv.lean}
(\srcid{92fd264bea303cb5a4f68c59d7a35bb0af8079c9}).
